% compendium/physical-commodity-markets/physical-markets.tex
% Entries migrated on 2026-10-16 from the legacy macro chapter (slice 37, 02_transmission_maps.tex), for crude-oil-products,
% metals-storage-delivery and agricultural-futures-delivery. Gliner (2014) is cited by chapter, section and PDF page of
% the ebook; his market shares and contract details are for 2011--2013. Dispositions in notes/migration-slices.md, slice 37.

\section{Crude oil, metals and agricultural markets}

\begin{qfdefinition}[Crude oil markets: grades, WTI and Brent, OPEC and spare capacity, chokepoints and pipelines, and the 3:2:1 crack spread, as of 2012--2013]\label{def:crude-oil-markets}
\qffield{Statement} Oil supplied 33 percent of world energy demand in 2012, coal 30 and natural gas 24 \citep[chap.~10, sec.\ ``Energy'', PDF p.~301]{s7-gliner}. Quality is measured by API gravity (higher means lighter; below 10 it sinks in water; 25 or less is ``heavy'') and sulphur content (sweet against sour), light sweet grades being the most valuable (pp.~301--302). OPEC has 12 members producing over 35 percent of world output and holding more than 70 percent of reserves; Saudi Arabia has the largest spare capacity, so ``the Call on OPEC'' usually means Saudi production rises (pp.~302--305, Table~10.4, OPEC-12 quota 30 million barrels a day as of September 2013). WTI, the U.S.\ benchmark, is delivered at Cushing, Oklahoma, and trades on NYMEX (CL), 1,000 barrels a contract; Brent, from North Sea fields, ``is used to price more than half of the world's internationally traded crude oil'' and trades on ICE (CO), 1,000 barrels; Brent's premium over WTI in Gliner's time reflected North Sea depletion, seaborne supply risk and pipeline constraints at Cushing (pp.~305--307). World consumption was 90 million barrels a day in 2012; the Strait of Hormuz carried 17 million barrels a day in 2011, ``about 20 percent of global traded oil and 35 percent of seaborne oil'', Malacca 14 million, Bab el-Mandeb 3 million and the Suez Canal 2 million; the Druzhba pipeline runs over 2,300 miles and carries up to 1.4 million barrels a day (pp.~308--310). The \emph{crack spread} is the refining margin; the 3:2:1 spread (three barrels of crude against two of gasoline and one of heating oil) is the most popular because gasoline output is almost double heating oil's, with product prices in cents a gallon converted at 42 gallons a barrel (sec.\ ``Crack Spread'', pp.~311--312).

\qffield{Notation} 3:2:1 crack per barrel $=\big[2\times(\text{RBOB}\times0.42)+1\times(\text{HO}\times0.42)-3\times\text{CL}\big]/3$, products in cents per gallon and crude in dollars per barrel; Gliner's worked example gives \$27.95 a barrel (p.~312), and crude at \$80, gasoline at \$2.40 and heating oil at \$2.70 a gallon give \$25.00 (slice-37 check script, section~14).

\qffield{Conditions} All shares, capacities, membership and contract details are Gliner's, as of 2011--2013. OPEC's membership, benchmark composition and trade flows may have changed since; none of this was checked against OPEC, EIA or exchange sources (backlog).

\qffield{Practice} The crack spread is a margin before operating costs; refiners' actual yields and product slates differ from 3:2:1. Political risk in producers and chokepoint disruptions move Brent first, in Gliner's account.

\qffield{Pitfalls} \emph{Errors in Gliner and the legacy.} ``Kerosene (diesel, jet fuel)'' follows Gliner's ``kerosene is used for diesel'' (p.~301); kerosene is jet fuel's base and diesel is a separate distillate. Gliner calls Saudi Aramco ``private'' and lists it among the ``biggest IOCs'' (pp.~303--304); it is a state-owned national oil company. ``Brent \dots\ Delivery at Sullom Voe'' repeats Gliner's ``deliverable to Sullom Voe'' (p.~306). ICE now describes its Brent Crude futures as ``a deliverable contract based on EFP delivery with an option to cash settle against the ICE Brent Index price for the last trading day'', 1,000 barrels, trading to the last business day of the second month preceding the contract month \citep[contract specifications, retrieved 2026-10-16]{ice-brent-spec}: delivery is by exchange for physical, not to a terminal, and there is a cash-settlement option. NYMEX WTI terms were not read (CME refused automated access). ``Historically WTI traded at premium (higher grade)'' is the legacy's; Gliner says only that WTI is the higher grade.

\qffield{See also} \hyperref[def:natural-gas-market]{natural gas}; \hyperref[def:contango-backwardation-stocks]{contango and backwardation}; \hyperref[def:commodity-indices-risk-return]{commodity indices}.

\qffield{Further reading} The pathway's readings on benchmarks, price reporting and the negative WTI settlement of April 2020.

\qffield{Index} crude oil; API gravity; sweet and sour crude; WTI; Brent; Cushing; OPEC; spare capacity; call on OPEC; Saudi Aramco; national oil companies; Strait of Hormuz; Strait of Malacca; chokepoints; Druzhba pipeline; Keystone; crack spread; 3:2:1 crack spread; refining margin

\qffield{Status} Verification: partly checked 2026-10-16 (Gliner's text, chap.~10, PDF pp.~301--312; the crack-spread arithmetic in the slice's check script; slice-37 follow-up, 2026-10-16: ICE Brent Crude futures specifications read; OPEC, EIA and CME pages not retrievable). Stability: draft. Currency: market data as of 2011--2013.
\end{qfdefinition}

\begin{qfdefinition}[Metals markets for macro: gold and real rates, silver, platinum, copper and aluminium, with their main producers, uses and contracts as of 2013]\label{def:metals-markets}
\qffield{Statement} \emph{Gold} (XAU) \citep[chap.~10, sec.\ ``Precious Metals'', PDF pp.~319--321]{s7-gliner}: jewellery ``more than 40 percent of demand'' and investment the second driver; central banks held 31,920 tonnes in October 2013, the United States 8,134 tonnes, 73 percent of its reserves (Table~10.6); ``China has overtaken'' South Africa ``as the leading producer'', and ``China, the United States, Australia, and South Africa collectively account for more than 40 percent of global gold production'' (p.~320); India, the largest consumer, used over 500 tonnes for jewellery and 300 for investment in 2011, ``of which 50 percent is used for weddings''; gold is preferred to fiat money when real interest rates are negative (Barsky and Summers's Gibson's paradox); GLD trades at a tenth of an ounce and COMEX gold (GC) is 100 troy ounces. \emph{Silver} (XAG, pp.~321--323): ``the highest electrical conductivity of all metals''; most demand industrial; ``Mexico and China account for 47 percent of global production'' (p.~322, on his Table~10.9); COMEX silver (SI) 5,000 troy ounces with ``triple 9s'' deliverable; SLV trades at one ounce. \emph{Platinum} (pp.~323--324): South Africa holds 80 percent of reserves; catalytic converters 38 percent and jewellery 31 percent of demand; COMEX platinum (PL) 50 troy ounces; traders follow the platinum/gold ratio. \emph{Copper} (sec.\ ``Industrial Metals'', pp.~325--327): ``Dr.\ Copper'', a leading indicator; production led by Chile and Peru; China the second producer and largest importer, about 40 percent of imports; uses equipment manufacturing 52 percent, building construction 33 and infrastructure 15; LME copper in dollars a tonne, 25 tonnes a contract, and COMEX copper (HG) in cents a pound, 25,000 pounds. \emph{Aluminium} (pp.~327--328): the most abundant metal in the crust; China over 40 percent of both production and consumption and self-sufficient; energy-intensive to produce; LME aluminium 25 tonnes.

\qffield{Notation} Troy ounce 31.1 grams; LME 3-month rolling forward quoted in dollars per tonne.

\qffield{Conditions} Figures are Gliner's for 2011--2013, from the World Gold Council, GFMS, ICSG, USGS and the World Copper Bulletin as he cites them; production shares were checked against USGS in the slice-37 follow-up (Pitfalls), consumption shares and exchange terms were not (backlog).

\qffield{Practice} The negative relation between gold and real rates is an empirical tendency with long episodes of breakdown; the legacy's ``gold inversely correlated with real interest rates'' should carry that qualification.

\qffield{Pitfalls} Gliner writes that platinum ``is almost twice has [sic] heavy as gold'' with densities of 21 and 19 (p.~323); the ratio is about 1.11 (checks, section~14). He writes that Chile represents ``more than 40 percent of consumption'' of copper (p.~325), which cannot be right for Chile, a producer and exporter; it probably means China, which he calls the largest importer in the same paragraph, but the text does not say so; the legacy's Chile at 27 percent of production and China at 43 percent of consumption were not matched to his table. ``India largest consumer (50\% used for weddings)'' compresses his statement about India's 2011 demand. ``China overtook South Africa as largest producer'' (legacy) is stated by Gliner (p.~320). \emph{Production shares checked against the data} (slice-37 follow-up). \emph{Gold}: USGS's \emph{Mineral Commodity Summaries 2013} gives 2011 and 2012 (estimated) mine output of 362 and 370 tonnes for China, 234 and 230 for the United States, 258 and 250 for Australia and 181 and 170 for South Africa, against world totals of 2,660 and 2,700 tonnes: 38.9 and 37.8 percent for the four \citep[gold, p.~2]{usgs-mcs-2013}, a little under Gliner's ``more than 40 percent'', whose source for that sentence is not stated (his Table~10.7 cites the World Gold Council, GFMS and the LBMA). \emph{Silver}: Gliner's Table~10.9 (World Silver Survey 2013 and GFMS) is internally inconsistent. Its ounces (Mexico 162.2, China 117.0 million) sum, with Other 120.0, to its stated total of 787.0 million, on which Mexico and China are $279.2/787.0=35.5$ percent and the top five 62.6 percent; but its printed shares (27.3 and 19.7 percent, and 120.0 as 20.2 percent) imply a denominator near 594 and add to 132.5 percent; his ``47 percent'' and ``83 percent'' for the top five follow the printed shares. USGS's 2013 summary gives 2012 estimates of 4,250 tonnes for Mexico and 3,800 for China against a world total of 24,000, 33.5 percent \citep[silver, p.~2]{usgs-mcs-2013}, close to the 35.5 percent the table's own ounces imply (slice-37 follow-up check script, section~4). The entry quotes Gliner as written and records the discrepancy rather than repairing it; the World Silver Survey itself was not read. \emph{Current shares} (USGS, \emph{Mineral Commodity Summaries 2025}, 2024 estimates): gold, China 380 of 3,300 tonnes, and China, Russia, Australia, Canada and the United States together 41 percent; silver, Mexico 6,300 and China 3,300 of 25,000 tonnes (38 percent); copper mine output led by Chile (5,300 of 23,000 thousand tonnes) and refined output by China (12,000 of 27,000) \citep[gold, silver and copper, pp.~1--2]{usgs-mcs-2025}; these are production shares, not the consumption shares the legacy quoted.

\qffield{See also} \hyperref[def:contango-backwardation-stocks]{contango and backwardation}; \hyperref[def:tips-breakevens]{real yields}; \hyperref[def:commodity-indices-risk-return]{commodity indices}.

\qffield{Further reading} The pathway's readings on metal inventories, exchange delivery, loco settlement and leases.

\qffield{Index} gold; central bank gold; real interest rates and gold; GLD; COMEX gold; silver; COMEX silver; platinum; platinum/gold ratio; copper; Dr.\ Copper; LME copper; COMEX copper; aluminium; LME aluminium; base metals; precious metals

\qffield{Status} Verification: partly checked 2026-10-16 (Gliner's text, chap.~10, PDF pp.~319--328; the density ratio in the slice's check script; slice-37 follow-up, 2026-10-16: Gliner pp.~320 and~322 reread and Table~10.9 recomputed; USGS \emph{Mineral Commodity Summaries} 2013 (gold, silver) and 2025 (gold, silver, copper) read). Stability: draft. Currency: market data as of 2011--2013.
\end{qfdefinition}

\begin{qfdefinition}[Agricultural markets for macro: corn, wheat, soybeans, cotton and coffee, crop calendars and sensitive stages, CBOT and ICE contracts, ethanol demand, and weather with the livestock cascade]\label{def:agricultural-markets}
\qffield{Statement} Gliner: agricultural supply depends on weather and planting, and ``some of the best agriculture traders have full-time meteorologists on staff''; crop calendars show the sensitive periods \citep[chap.~10, sec.\ ``Agriculture'', PDF p.~330]{s7-gliner}. \emph{Corn} (pp.~330--335): the United States produces over 35 percent and consumes over 30 percent of the world's corn and is the largest exporter; Iowa, Illinois and Nebraska produce 47 percent of U.S.\ corn; uses are feed, ethanol (27 percent of U.S.\ corn demand in 2011, after the Energy Policy Act of 2005) and food; Chinese feed imports rose from nothing in 2007 to over 5 percent of world imports; silking from June is the critical period; CBOT corn (C) is quoted in cents a bushel, 5,000 bushels a contract, 56 pounds a bushel, for March, May, July, September and December. \emph{Wheat} (pp.~335--338): classified by season, gluten and colour; the United States is the fourth producer but a leading exporter; a close substitute for corn; heading is the sensitive stage; CBOT wheat (W), 5,000 bushels of 60 pounds. \emph{Soybeans} (pp.~338--341): the United States the largest producer and second exporter (31 and 37 percent), Brazil the largest exporter, China 63 percent of imports; more than 90 percent genetically engineered; flowering is the vulnerable stage; CBOT soybeans (S), 5,000 bushels of 60 pounds, for January, March, May, July, August, September and November. \emph{Cotton} (pp.~341--344): character, grade and staple; China, India and the United States 65 percent of production in 2013, China the largest producer and consumer, the United States the largest exporter with almost 30 percent; ICE cotton (CT) in cents a pound, 50,000 pounds, roughly 104 bales of 480 pounds. \emph{Coffee} (pp.~344--348): arabica about 60 percent of production, robusta the rest; Brazil the largest producer (37.2 percent in 2012--13) and Vietnam second (16.6 percent), in bags of 60 kilograms (Table~10.21); the United States and the European Union almost half of demand; ICE arabica (KC). \emph{Weather and livestock} (p.~333): when grain prices rise sharply, feeding livestock becomes expensive and farmers slaughter early, so near-term livestock prices fall while later prices rise.

\qffield{Notation} Futures month codes F (Jan), H (Mar), K (May), N (Jul), Q (Aug), U (Sep), V (Oct), X (Nov), Z (Dec); one bushel of corn 56 pounds, of wheat and soybeans 60.

\qffield{Conditions} Shares, contract terms and delivery months are Gliner's for 2012--2013 (USDA tables as he reproduces them); they were not checked against current CBOT and ICE rulebooks or USDA data (backlog). A corn contract at 500 cents is \$25,000 notional (checks, section~14).

\qffield{Practice} Translate a view on a crop's balance sheet (\hyperref[def:contango-backwardation-stocks]{ending stocks}) into a specific contract month, grade and delivery location; the basis between local cash and futures is the pathway's subject.

\qffield{Pitfalls} \emph{Legacy additions not in the passages read}: El Ni\~no and La Ni\~na, NOAA drought monitors, ``demand relatively inelastic'', Russian wheat export bans, coffee's biennial cycle and 20-year tree life; plausible, but unchecked here. Two are now checked against ICE (slice-37 follow-up): Coffee ``C'' (KC) is 37,500 pounds, quoted in cents a pound, for March, May, July, September and December, physically delivered; Cotton No.~2 (CT) is 50,000 pounds net weight, for March, May, July, October and December, matching the legacy's H, K, N, V, Z \citep[contract specifications, retrieved 2026-10-16]{ice-softs-specs}. CBOT grain terms were not read (CME refused automated access). The legacy's cotton contract size ``50,00 0 lbs'' is a typo for 50,000.

\qffield{See also} \hyperref[def:contango-backwardation-stocks]{contango and ending stocks}; \hyperref[def:futures-hedge-basis]{futures hedges and basis}; \hyperref[def:commodity-indices-risk-return]{commodity indices}.

\qffield{Further reading} WASDE \citep{physical-usda-wasde-2026} and the pathway's readings on grain delivery.

\qffield{Index} agriculture; corn; wheat; soybeans; cotton; coffee; arabica; robusta; crop calendar; silking; heading; flowering; CBOT grain contracts; ICE softs; bushel; ethanol demand; weather risk; livestock cascade; genetically engineered crops

\qffield{Status} Verification: partly checked 2026-10-16 (Gliner's text, chap.~10, PDF pp.~330--348; slice-37 follow-up, 2026-10-16: ICE coffee and cotton specifications read). Stability: draft. Currency: market data as of 2012--2013.
\end{qfdefinition}
